% -*- coding: utf-8 -*-
% !TEX root = ../main.tex

\subsection{Conclusion}
\label{ssec:conclusion_conclusion}

Ce travail avait pour objectif d'étudier le compromis entre qualité de
classification, mémoire et calcul lors de la compression d'une tête de
classification de sentiments. Un pipeline reproductible a été construit autour
d'un tokenizer BPE, d'un backbone CBOW figé et de deux têtes neuronales. Cette
séparation a permis de modifier uniquement le classifieur tout en conservant
exactement les mêmes représentations d'entrée. Deux critères de pruning et un
témoin aléatoire ont été comparés : OBD, fondé sur une approximation diagonale
de la Hessienne de la loss ; MINT, fondé sur une estimation géométrique de
l'information mutuelle conditionnelle entre activations ; et un masque tiré
uniformément sans utiliser les poids, les activations ou les labels.

Les expériences montrent qu'une cible proche de $50\,\%$ de poids supprimés est
compatible avec le maintien de la performance finale. Sur cinq graines,
H1/OBD-M3 atteint un F1 de $0{,}7512\pm0{,}0048$ contre
$0{,}7456\pm0{,}0048$ pour M1. Sur H2, MINT-M3 atteint
$0{,}7430\pm0{,}0101$ contre $0{,}7321\pm0{,}0082$. Les intervalles des
différences appariées contiennent néanmoins zéro et de faibles valeurs
négatives. Avec cinq graines,
les résultats sont donc compatibles avec une légère baisse, une absence d'effet
ou une hausse du F1 ; ils ne démontrent formellement ni une amélioration, ni une
non-dégradation, ni un effet de régularisation dû au seul masque. MINT se
distingue par une forte fragilité immédiate : avant
réentraînement, son F1 moyen tombe à $0{,}2028\pm0{,}0938$. OBD produit un
masque directement plus stable et son scoring est environ 163 fois plus rapide
que celui de MINT sur H2.

Le témoin aléatoire apporte une information utile. Sur H1, OBD-M3 dépasse
RAND-M3 de $0{,}0177$ de F1 moyen et l'intervalle apparié à $95\,\%$ ne contient
pas zéro. OBD choisit donc de meilleurs poids qu'un tirage au hasard sur cette
tête. Sur H2, RAND-M3 atteint $0{,}7323\pm0{,}0139$, contre
$0{,}7371\pm0{,}0068$ pour OBD-M3 et $0{,}7430\pm0{,}0101$ pour MINT-M3. Les
intervalles des deux écarts contiennent zéro. Ces résultats ne permettent donc
pas de dire qu'OBD ou MINT choisit un meilleur sous-réseau sur H2.

La conclusion la plus importante concerne la différence entre sparsité et
compression déployable. Les masques divisent par deux le nombre idéal de
paramètres et de MAC de la tête, mais ils ne changent ni la forme des tenseurs,
ni la taille de l'artefact dense, ni sa latence. La petite architecture M4,
entraînée directement, obtient un F1 moyen de $0{,}7345\pm0{,}0100$ sur H2
tout en réduisant
réellement la taille de l'artefact de $34{,}95\,\%$ et la latence bout en bout
de $12{,}83\,\%$. Dans l'environnement actuel, M4 offre donc le compromis le
plus immédiatement exploitable. OBD et MINT identifient des sous-réseaux
performants, mais leur potentiel ne deviendra un gain matériel qu'après une
conversion vers une représentation ou une architecture qui exploite les zéros.

\subsection{Retours d'expérience}
\label{ssec:conclusion_retours}

Ce projet a d'abord mis en évidence l'importance de formaliser précisément ce
que signifie \emph{comprimer}. Compter les poids non nuls est utile pour
comparer les masques, mais ne renseigne pas à lui seul sur la mémoire physique
ou le temps d'exécution. L'ajout de deux vues --- tête seule et pipeline complet
--- puis de mesures de latence et de stockage a évité de confondre un gain
algorithmique idéal avec un gain effectivement obtenu.

La mise en œuvre d'OBD a permis d'approfondir l'autodifférentiation d'ordre
deux, les produits Hessienne--vecteur et l'estimation stochastique de la
diagonale. Celle de MINT a demandé de relier une définition informationnelle à
des objets calculables : activations standardisées, bootstrap conditionnel par
plus proche voisin et statistique de Friedman--Rafsky sur un arbre couvrant
minimal. Leur comparaison a aussi montré qu'une méthode plus riche
conceptuellement n'est pas automatiquement plus efficace : le coût de scoring,
la sensibilité au masque et le besoin de réentraînement font partie intégrante
du bilan. Le témoin aléatoire permet aussi de séparer la qualité du score de la
capacité du réseau à récupérer après le pruning.

Enfin, la construction d'un pipeline unique, appuyé sur une configuration
centralisée et des artefacts archivés, a renforcé la reproductibilité du
travail. La principale difficulté n'a pas seulement été d'obtenir une bonne
métrique, mais de garantir que les modèles comparés utilisent les mêmes données,
le même backbone, la même cible de capacité et des définitions cohérentes de la
complexité.

\subsection{Perspectives}
\label{ssec:conclusion_perspectives}

La première perspective est statistique. Les cinq graines actuelles pourraient
être étendues à un effectif plus important, puis complétées par plusieurs splits
ou une validation croisée afin de ne plus conditionner l'analyse à un unique
jeu de test et à un unique W1. Le batch de scoring OBD devra être tiré
explicitement, et les scores MINT devront être moyennés sur plusieurs partitions
afin d'estimer leur stabilité interne. Une correction des comparaisons multiples
et une calibration du seuil de classification sur la validation permettraient
également de séparer la perte de classement d'un simple déplacement des
probabilités après pruning.

La deuxième perspective consiste à transformer la sparsité en gain physique.
Trois voies peuvent être comparées : stockage sparse des matrices, noyaux
d'inférence spécialisés et pruning structuré supprimant des groupes complets
d'unités. Cette dernière solution permettrait de reconstruire une tête compacte
à partir du masque, puis de la comparer à M4 à dimensions et conditions
d'entraînement contrôlées. Les mesures devraient alors couvrir plusieurs tailles
de batch, le CPU et le GPU, la mémoire maximale et, si possible, l'énergie
consommée.

La troisième perspective concerne les critères de scoring. Le calcul MINT
actuel répète une estimation géométrique pour chaque couple unité--groupe. La
mise en cache des distances, la parallélisation des blocs ou la réduction de
dimension du conditionnement pourraient diminuer ses $230$ secondes de
scoring. Plus généralement, OBD et MINT pourraient être comparés à des scores
fondés sur la magnitude, sur une approximation du premier ordre par le gradient,
sur l'information de Fisher ou sur une approximation tenant compte des termes
hors diagonale de la Hessienne. Des critères directement définis sur des
neurones ou des blocs seraient également cohérents avec l'objectif d'obtenir
une architecture physiquement compacte. La comparaison devrait conserver la
même sparsité cible, le même budget de réentraînement et les mêmes graines, à
l'image du contrôle aléatoire désormais inclus.

Une autre limite importante se situe avant la tête de classification. Le
pipeline représente actuellement un texte par la moyenne des embeddings de ses
unités BPE. Cette opération est simple et peu coûteuse, mais elle est invariante
à leur permutation : les phrases « le chien mord l'homme » et « l'homme mord le
chien » produiraient le même vecteur si elles contenaient exactement les mêmes
unités. Or leur interprétation dépend de l'ordre. Il serait donc pertinent de
comparer la moyenne à une agrégation attentive munie d'informations de position,
à un réseau convolutif capturant des motifs locaux, ou à un encodeur récurrent
de type GRU ou LSTM. Ces méthodes remplaceraient le sac d'embeddings actuel par
une représentation sensible à la séquence.

Cette évolution permettrait aussi d'étudier des têtes plus avancées, notamment
un encodeur Transformer. Pour exploiter réellement l'ordre, celui-ci ne devrait
pas recevoir le vecteur déjà moyenné, mais la séquence des embeddings W1,
complétée par un encodage positionnel, puis produire une représentation globale
par attention ou par un jeton dédié à la classification. W1 pourrait rester
figé, conformément au périmètre du projet, tandis que le Transformer ferait
partie de la tête à comprimer. Une telle architecture augmenterait cependant
le coût initial du classifieur ; son gain de qualité devrait donc être rapporté
à sa mémoire, à ses opérations et à sa latence réelles.

Le cadre expérimental pourrait enfin être étendu au multiclasse. Le problème
actuel possède déjà deux classes, positive et négative, car les exemples neutres
ont été retirés. Une première extension naturelle consisterait à les réintroduire
pour obtenir $K=3$, puis à considérer des jeux comportant davantage de classes.
La tête produirait alors un logit par classe et serait entraînée avec une
cross-entropy multiclasse. Le macro-F1, les rappels par classe et la matrice de
confusion deviendraient essentiels pour vérifier qu'un score de pruning ne
préserve pas seulement les classes majoritaires.

Une direction distincte consisterait à remplacer W1 par des représentations
contextuelles préentraînées. Le backbone devrait toujours rester identique entre
les variantes afin de conserver l'interprétabilité de la comparaison. Cette
extension permettrait de vérifier si les conclusions obtenues avec CBOW ---
redondance importante de la tête, nécessité de cicatrisation et écart entre
sparsité théorique et vitesse réelle --- se maintiennent lorsque les
représentations d'entrée sont elles-mêmes sensibles au contexte.
